\appendix
\chapter{HỒ SƠ TÁI LẬP VÀ HƯỚNG PHÁT TRIỂN}
\section{Tái lập đánh giá VALID}
Quy trình bắt đầu bằng kiểm tra ROD Final A4 có 18.151 ảnh TRAIN, 3.318 ảnh VALID, 1.611 ảnh TEST và hash \texttt{data.yaml} khớp khóa. Tiếp theo kiểm SHA256 ba checkpoint. Không dùng checkpoint cùng tên nhưng khác hash, vì chỉ một thay đổi trọng số cũng tạo thí nghiệm khác.

E3/E4 nhận ảnh ở \texttt{imgsz=640}, thu prediction tại confidence 0,01, NMS 0,7 và \texttt{max\_det=300}. E6 dựng Faster R-CNN ResNet-50 FPN V2, thay head 26 lớp, nạp state dict strict và ánh xạ label 1--25 về 0--24. Prediction được lưu theo image ID để tính lại threshold mà không chạy model lần nữa.

Evaluator kiểm không thiếu/trùng ID, rồi tính AP từng lớp tại mười ngưỡng IoU. Prediction được sắp score giảm dần và ghép greedy với GT cùng lớp chưa dùng. Đường Precision--Recall nội suy 101 điểm. mAP50 trung bình 25 lớp tại IoU 0,5; mAP50--95 trung bình thêm qua mười ngưỡng IoU. Precision/Recall vận hành dùng ngưỡng F2 khóa trước và IoU 0,5.

Khi công bố kết quả cần lưu JSON tổng, CSV theo lớp, hash cache prediction, checkpoint, data YAML và protocol. Nếu chỉ lưu ảnh chụp hoặc một số mAP, không thể xác định dữ liệu, ngưỡng hay code đã dùng. Báo cáo ưu tiên tệp máy đọc vì lý do này.

\section{Đánh giá TEST cuối}
Trước khi mở TEST, kiến trúc, checkpoint, confidence threshold, NMS, max detection, mapping lớp và cách tính metric phải đóng băng. \texttt{evaluation\_lock.json} là biên bản kỹ thuật. Script phải từ chối nếu hash không khớp.

Ba model sau đó chạy trên đúng 1.611 ảnh TEST. Ngưỡng E3/E4/E6 vẫn là 0,22/0,19/0,59; không điều chỉnh theo TEST. Cần báo Precision, Recall, F2, mAP50, mAP50--95, TP/FP/FN và bảng 25 lớp. Nếu model thắng VALID nhưng không thắng TEST, phải giữ cả hai kết quả và phân tích overfit lựa chọn.

Một lượt TEST không cung cấp biến thiên theo seed. Có thể bootstrap theo ảnh để ước lượng độ bất định chênh lệch metric, nhưng phải giữ nguyên prediction và lấy mẫu ảnh với hoàn lại. Nếu huấn luyện nhiều seed, từng seed cần checkpoint và log riêng; không chọn seed tốt nhất bằng TEST.

\section{Phân tích lỗi định tính}
Mỗi FN nên được gắn thuộc tính: vật nhỏ, che khuất, cắt biên, ánh sáng yếu, motion blur, nền phức tạp, lớp gần nghĩa hoặc nhãn nghi vấn. FP được phân thành nhầm lớp, trùng hộp, nhận nền thành vật thể hoặc phát hiện vật thể thật nhưng thiếu nhãn. Hộp đúng lớp nhưng IoU thấp được tách khỏi lỗi phân loại. Phân loại lỗi cho biết nên cải thiện dữ liệu, augmentation hay kiến trúc.

Mẫu lỗi phải lấy theo quy tắc thay vì chọn ví dụ gây ấn tượng. Có thể lấy ngẫu nhiên theo lớp và loại lỗi, đồng thời báo cỡ mẫu. Với Person, Car và Pedestrian crosswalk có nhiều FN tuyệt đối, cần xem thêm tỷ lệ trên số GT. Với lớp hiếm, vài lỗi có thể làm Recall thay đổi lớn.

Sai nhãn phải ghi vào manifest review và không sửa trực tiếp trong TEST sau khi xem prediction. Nếu xác nhận TEST có lỗi annotation, nên công bố cả kết quả trên bản đóng băng và phiên bản mới cho nghiên cứu sau; không âm thầm sửa để cải thiện model.

\section{Đánh giá trên phần cứng mục tiêu}
Benchmark CPU local 30 ảnh, batch 1 cho phép so tương đối nhưng chưa đại diện smartphone hoặc GPU. Đánh giá triển khai phải cố định thiết bị, driver/framework, chế độ nguồn, kích thước ảnh, số luồng và warmup; báo mean, median, p95, request đầu, bộ nhớ đỉnh và nhiệt độ nếu chạy dài.

Nếu chuyển sang video, FPS ảnh tĩnh chưa đủ. Hệ thống cần đo camera capture, resize, inference, tracking, vẽ, âm thanh và hàng đợi; khi xử lý chậm hơn camera phải quy định bỏ frame. Độ trễ cảnh báo quan trọng hơn throughput trung bình vì frame chậm có thể tạo khoảng im lặng nguy hiểm.

Tối ưu có thể gồm ONNX, TensorRT, TorchScript, lượng tử hóa hoặc giảm kích thước ảnh. Mỗi thay đổi phải đánh giá lại accuracy vì tăng tốc có thể ảnh hưởng hộp nhỏ. E3 có lợi thế dung lượng; E4 cân bằng; E6 cần tăng tốc mạnh nếu nhắm thời gian thực trên CPU. Không suy đoán mức tăng từ tài liệu nhà cung cấp mà phải đo checkpoint ROD.

\section{Giới hạn liên quan người dùng khiếm thị}
Detector ảnh tĩnh chỉ trả tên và vị trí pixel. Nó không biết vật thể cách bao xa, có nằm trên quỹ đạo hay đang tiến lại gần. Cảnh báo mọi detection có thể gây quá tải; bỏ qua detection có thể gây rủi ro. Thiết kế cảnh báo cần nghiên cứu người dùng, ưu tiên theo nguy cơ, cooldown, tracking và thử nghiệm có giám sát.

Âm thanh phải ngắn, dễ hiểu và không che âm thanh môi trường. Phản hồi rung cần quy ước nhất quán. Các quyết định này không thể suy từ mAP. Trước thử nghiệm người dùng cần đánh giá an toàn, đồng thuận, quyền riêng tư hình ảnh và cơ chế dừng khẩn cấp. Ứng dụng hiện chỉ là demo nghiên cứu ảnh tải lên.

\section{Ánh xạ 25 lớp}
\begin{longtable}{rll}
\toprule ROD ID & Torchvision ID & Tên lớp\\\midrule\endfirsthead
\toprule ROD ID & Torchvision ID & Tên lớp\\\midrule\endhead
0&1&Bike\\1&2&Building\\2&3&Car\\3&4&Person\\4&5&Stairs\\
5&6&Traffic sign\\6&7&Electrical Pole\\7&8&Road\\8&9&Motorcycle\\9&10&Dustbin\\
10&11&Dog\\11&12&Manhole\\12&13&Tree\\13&14&Guard rail\\14&15&Pedestrian crosswalk\\
15&16&Truck\\16&17&Bus\\17&18&Bench\\18&19&Traffic Cone\\19&20&Fire hydrant\\
20&21&Teraffic Barrel\\21&22&Plant Pot\\22&23&Electrical Box\\23&24&Chair\\24&25&Bicycle Rack\\\bottomrule
\end{longtable}

Torchvision ID 0 là background và không ánh xạ thành lớp ROD. YOLO dùng trực tiếp ROD ID 0--24. Bảng phải được kiểm khi thay dataset hoặc checkpoint; sai một đơn vị có thể tạo hộp đúng nhưng tên lớp sai cho toàn bộ ứng dụng.
